% Part: first-order-logic
% Chapter: beyond
% Section: modal-logics

\documentclass[../../../include/open-logic-section]{subfiles}

\begin{document}

% SEGMENT OLP-0180-S01

\olfileid{fol}{byd}{mod}

\olsection{வாய்ப்புநிலைத் தருக்கங்கள்}

பின்வரும் நிபந்தனை வாக்கிய எடுத்துக்காட்டைக் கருதுக:
\begin{quote}
  அந்த அறையில் ஜெரமி தனியாக இருந்தால், அவர் குடிபோதையில் ஆடையின்றி
  நாற்காலிகளின் மீது ஆடிக்கொண்டிருக்கிறார்.
\end{quote}
இது பொருள்சார் நோக்கில் மெய்யாகவும், இருந்தும் தவறான எண்ணத்தை
ஏற்படுத்துவதாகவும் இருக்கக்கூடிய ஒரு நிபந்தனைக் கூற்றுக்கு எடுத்துக்காட்டு.
முன்னுறுப்பு பொய்யாகவோ பின்னுறுப்பு மெய்யாகவோ இருப்பதை மட்டும் கடந்து,
முன்னுறுப்புக்கும் முடிவுக்கும் இடையே வலுவான தொடர்பு இருப்பதாக இது
உணர்த்துவதுபோல் தெரிகிறது. அதாவது, இந்தக் குறிப்பிட்ட உலகத்தில்
(ஜெரமி அந்த அறையில் தனியாக இல்லாததால் அற்பமான முறையில் மெய்யாக
இருக்கக்கூடிய நிலையில்) மட்டும் இக்கூற்று மெய்யல்ல; முன்னுறுப்பு
மெய்யாக \emph{இருந்திருந்தால்}, முடிவும் மெய்யாக
\emph{இருந்திருக்கும்} என்பதையும் சொற்கள் உணர்த்துகின்றன. வேறுவிதமாகச்
சொன்னால், இவ்வுலகத்தில் மட்டுமல்லாமல் எந்த ``சாத்தியமான'' உலகத்திலும்
இக்கூற்று மெய் என்றும், அல்லது இக்குறிப்பிட்ட உலகத்தில் மட்டும் மெய்
என்பதற்கு மாறாக அது \emph{இன்றியமையாமுறையில்} மெய் என்றும் இக்கூற்றைப்
பொருள்கொள்ளலாம்.

இவ்வகை இன்றியமையாமைக்குப் பொருள் வழங்க வாய்ப்புநிலைத் தருக்கம்
வடிவமைக்கப்பட்டது. வழக்கமான கூற்றுக்கணிப்புத் தருக்கத்துடன் பெட்டி
செயற்குறியைச் சேர்ப்பதன் மூலம் வாய்ப்புநிலைக் கூற்றுக்கணிப்புத் தருக்கம்
கிடைக்கிறது; அதாவது $!A$ ஓர் !!a{formula} எனில், $\Box !A$ உம்
அவ்வாறே. உள்ளுணர்வாக, $\Box !A$ ஆனது $!A$
\emph{இன்றியமையாமுறையில்} மெய், அல்லது சாத்தியமான எந்த உலகத்திலும்
மெய் என்று கூறுகிறது. $\Diamond !A$ என்பது வழக்கமாக
$\lnot \Box \lnot !A$ என்பதன் சுருக்கமாகக் கொள்ளப்பட்டு, $!A$
\emph{சாத்தியமாக} மெய் என்று கூறுவதாகப் படிக்கப்படுகிறது. வாய்ப்புநிலையைப்
பயனிலைக் கணிப்புத் தருக்கத்துடனும் சேர்க்கலாம்.

வாய்ப்புநிலைத் தருக்கத்திற்குப் பொருண்மையியலை வழங்க கிரிப்கே
!!{structure}s ஐப் பயன்படுத்தலாம்; உண்மையில், வாய்ப்புநிலைத்
தருக்கத்தையே மனத்தில் கொண்டு கிரிப்கே இப்பொருண்மையியலை முதலில்
வடிவமைத்தார். பகுதி வரிசைகளுக்கு மட்டும் கட்டுப்படுத்துவதற்குப் பதிலாக,
பொதுவாக ``சாத்தியமான உலகங்களின்'' கணம் $P$ ஒன்றும், உலகங்களுக்கிடையிலான
இரும ``அணுகுறவு'' $\Atom{R}{x,y}$ ஒன்றும் இருக்கும். உள்ளுணர்வாக,
$\Atom{R}{p,q}$ என்பது உலகம்~$q$ ஆனது~$p$ உடன் ஒத்திசைவானது என்று
கூறுகிறது; அதாவது நாம் உலகம்~$p$ இல் ``இருந்தால்'', உலகம்~$q$ போலவும்
இருந்திருக்கக்கூடிய சாத்தியத்தை நாம் கருத வேண்டும்.

``விரிநிலை'' தருக்கத்திற்கு மாறாக வாய்ப்புநிலைத் தருக்கம் சிலவேளைகளில்
``கருத்துநிலை'' தருக்கம் என்று அழைக்கப்படுகிறது. செவ்வியல் தருக்கம்
போன்ற விரிநிலைத் தருக்கத்திற்கான நோக்கம்கொண்ட பொருண்மையியல்,
``நடப்பு'' உலகமான ஒரே உலகத்தை மட்டுமே சுட்டும்; ஆனால்
``கருத்துநிலை'' தருக்கத்தின் பொருண்மையியல் இன்னும் விரிவான
மெய்ப்பொருளியலைச் சார்ந்துள்ளது. இன்றியமையாமையை !!{structure}ப்பதுடன்,
பயன்பாட்டுக்கு ஏற்ப $\Box$ மற்றும் $\Diamond$ ஐ மறுபொருள்கொண்டு,
பிற மொழியியல் கட்டமைப்புகளை !!{structure}க்கவும் வாய்ப்புநிலையைப்
பயன்படுத்தலாம். எடுத்துக்காட்டாக:
\begin{enumerate}
\item நிரூபிக்கத்தக்கமைத் தருக்கத்தில், $\Box !A$ என்பது
  ``$!A$ நிரூபிக்கத்தக்கது'' என்றும், $\Diamond !A$ என்பது
  ``$!A$ முரணின்மையுடையது'' என்றும் படிக்கப்படுகின்றன.
\item அறிவுநிலைத் தருக்கத்தில், $\Box !A$ என்பதை
  ``நான் $!A$ ஐ அறிவேன்'' அல்லது ``நான் $!A$ ஐ நம்புகிறேன்''
  என்று படிக்கலாம்.
\item காலநிலைத் தருக்கத்தில், $\Box !A$ என்பதை ``$!A$ எப்போதும்
  மெய்'' என்றும், $\Diamond !A$ என்பதை ``$!A$ சிலவேளைகளில் மெய்''
  என்றும் படிக்கலாம்.
\end{enumerate}

வாய்ப்புநிலையைக் கையாளும் விதிகளையும் அடிகோள்களையும் தருக்கத்துடன்
சேர்க்க விரும்புகிறோம். எடுத்துக்காட்டாக, வழக்கமான கூற்றுக்கணிப்புத்
தருக்க அடிகோள்கள் மற்றும் விதிகளுடன் பின்வரும் அடிகோள்களையும்
$\Log{S4}$ அமைப்பு கொண்டுள்ளது:
\begin{align*}
& \Box (!A \lif !B) \lif (\Box !A \lif \Box !B)\\
& \Box !A \lif !A \\
& \Box !A \lif \Box \Box !A
\intertext{அத்துடன் ``$!A$ இலிருந்து $\Box !A$ என்று முடிவுகொள்க''
  என்ற விதியும் உண்டு. $\Log{S5}$ பின்வரும் அடிகோளைச் சேர்க்கிறது:}
& \Diamond !A \lif \Box \Diamond !A
\end{align*}
இவ்வடிகோள்களின் மாற்றுவடிவங்கள் வெவ்வேறு பயன்பாடுகளுக்கு ஏற்றவையாக
இருக்கலாம்; எடுத்துக்காட்டாக, S5 என்பது தருக்க இன்றியமையாமைக்
கருத்துருவின் தன்மையைப் பொதுவாகக் குறிப்பதாகக் கொள்ளப்படுகிறது.
கிரிப்கே !!{structure}s இலுள்ள அணுகுறவை இயல்பான முறைகளில்
கட்டுப்படுத்துவதன் மூலம், !!{derivation} அமைப்பு சரித்தன்மையுடனும்
நிறைவுடனும் இருக்கும் பொருண்மையியலைப் பொதுவாகக் கண்டறிய முடிவது இதன்
சிறப்பு. எடுத்துக்காட்டாக, அணுகுறவு தற்பின்னோக்கியதாகவும் கடப்புநிலை
உடையதாகவும் இருக்கும் கிரிப்கே !!{structure}s இன் வகுப்பிற்கு
$\Log{S4}$ ஒத்துள்ளது. அணுகுறவு {\em பொதுமையானதாக}, அதாவது ஒவ்வோர்
உலகத்தையும் மற்ற ஒவ்வோர் உலகத்திலிருந்தும் அணுகக்கூடியதாக, இருக்கும்
கிரிப்கே !!{structure}s இன் வகுப்பிற்கு $\Log{S5}$ ஒத்துள்ளது; எனவே
ஒவ்வோர் உலகத்திலும் $!A$ மெய்யானால், எனில் மட்டுமே, $\Box !A$
மெய்யாகிறது.

\end{document}
